% ===========================================================================
\chapter{Giving the model hands}\label{ch:tools}

A tool is something the agent can do for the model: list the files, read
one, search them, edit one, run a command. The model can't do any of
these itself. It asks, by name, with a JSON object of arguments; the agent
checks the request, does it, and sends back what happened as text.

For the model to ask for a tool, it has to know the tool exists. Every
request lists the tools, each with a name, a description, and a
\emph{schema}: a description, in JSON, of the arguments it takes. Here is
the one for the file reader we write in this chapter:

\begin{json}
{"name": "read_file",
 "description": "Read a text file in the workspace. Lines come numbered, ...",
 "input_schema": {"type": "object",
                  "properties": {"path": {"type": "string", "description": "The file, ..."},
                                 "start_line": {"type": "integer", "description": "..."},
                                 "end_line": {"type": "integer", "description": "..."}},
                  "required": ["path"],
                  "additionalProperties": false}}
\end{json}

The model reads all of this, every time, and decides from it alone which
tool to use and how. That makes the descriptions part of the program, as
much as the code is. A good description says not only what a tool does but
\emph{when} to use it, and anything about its input or output that would
otherwise be a surprise.

\section{The Tool class}

A tool's run gives back an outcome: text, and whether it is an error. Two
constructors, \code{ok} and \code{error}, say which:

\begin{newfile}{tools.nex}
-- tools.nex: what the agent can do in its workspace.
--
-- The model never touches anything itself. It asks for a tool by name,
-- with a JSON object of arguments; the agent checks the request, runs
-- the tool, and sends back what happened as text.

intern Messages
intern Workspace
intern Helpers
intern data/Json

-- What running a tool gave back: some text, and whether it failed.
class Tool_Outcome
create
  ok(text: String) do
    this.text := text
    is_error := false
  end

  error(text: String) do
    this.text := text
    is_error := true
  end

feature
  text: String
  is_error: Boolean
end

-- One parameter of a tool: its JSON type ("string" or "integer") and
-- what it means, for the model to read.
function parameter(json_type, description: String): Map[String, Any]
require
  known_type: json_type = "string" or json_type = "integer"
do
  result := {"type": json_type, "description": description}
end
\end{newfile}

\code{Tool} is deferred, like \code{Model}: each tool says its name, its
description, its parameters and how it runs. Each tool works in a
workspace, which it is given when it is made. A tool also says whether
running it can change anything; tools that can will need permission
(Chapter~\ref{ch:permission}). And it can describe one call in a few
words, such as \code{read_file loans.nex}, for printing.

\begin{append}{tools.nex}
deferred class Tool
create
  make(ws: Workspace) do
    workspace := ws
  end

feature
  workspace: Workspace    -- where it works

  -- The name the model calls it by.
  name(): String deferred

  -- What it does and when to use it. The model decides which tool to
  -- call from this alone, so it says when, not just what.
  description(): String deferred

  -- The parameters, by name.
  parameters(): Map[String, Any] deferred

  -- The names of the parameters that must be given.
  required(): Array[String] deferred

  -- Whether running it can change anything: files, or whatever a
  -- command does. Such tools need permission.
  changes_things(): Boolean do
    result := false
  end

  -- A short description of one call, for the transcript and for
  -- permission questions: "read_file loans.nex".
  describe(input: Json_Value): String deferred

  -- Run the tool on `input`, already checked by problem_with.
  run(input: Json_Value): Tool_Outcome
  require
    valid_input: problem_with(input) = ""
  deferred

  -- The tool as the API wants it described.
  definition(): Map[String, Any] do
    let schema: Map[String, Any] := {"type": "object", "properties": parameters(),
      "required": required(), "additionalProperties": false}
    result := {"name": name(), "description": description(), "input_schema": schema}
  end
\end{append}

\section{Checking the input}

\code{run}'s precondition says its input must already have been checked.
Models follow schemas well, but not perfectly: now and then an argument is
missing, misnamed, or a number arrives as a string. Rather than have every
tool check its own input, \code{problem_with} checks it against the tool's
parameters, the same ones the model was shown. Its answer is \code{""} if
all is well, and otherwise a sentence the model can act on. The agent
sends that sentence back as an error, and the model nearly always corrects
itself on its next try.

\begin{continued}{tools.nex}
  -- What is wrong with `input`, or "" if nothing is. The model usually
  -- follows the schema, but not always, and a tool must not trust it.
  problem_with(input: Json_Value): String do
    result := ""
    if not input.is_object then
      result := "the input must be a JSON object"
    else
      across required() as key do
        if result = "" and not input.has(key) then
          result := "missing parameter `" + key + "`"
        end
      end
      let params: Map[String, Any] := parameters()
      across input.keys() as key do
        if result = "" then
          if not params.contains_key(key) then
            result := "unknown parameter `" + key + "`"
          elseif type_of_parameter(key) = "string" and not input.get(key).is_string then
            result := "`" + key + "` must be a string"
          elseif type_of_parameter(key) = "integer" and not is_whole_number(input.get(key)) then
            result := "`" + key + "` must be a whole number"
          end
        end
      end
    end
  end
\end{continued}

Nex has no \code{return} statement, so a routine that wants to stop at the
first problem it finds does it with a test: once \code{result} holds a
problem, the later checks are skipped.

Tools read their arguments with two helpers that fall back to a default.
They don't assume the input is valid, because \code{describe} is used on
input that hasn't been checked (we want to print what the model asked for,
even when it asked wrongly):

\begin{continued}{tools.nex}
  -- The string argument `key`, or `default` if it was not given (or is
  -- not a string: describe is used on input that has not been checked).
  string_argument(input: Json_Value, key, default: String): String do
    result := default
    if input.is_object and input.has(key) and input.get(key).is_string then
      result := input.get(key).as_string
    end
  end

  -- The whole-number argument `key`, or `default` if it was not given.
  integer_argument(input: Json_Value, key: String, default: Integer): Integer do
    result := default
    if input.is_object and input.has(key) and is_whole_number(input.get(key)) then
      result := input.get(key).as_integer
    end
  end
\end{continued}

Every tool that takes a path checks it the same way, so that is a helper
too. The workspace's contracts would stop a bad path anyway, but with a
precondition violation, which is the end of the program. Checked here, a
bad path is just an error for the model, with a hint:

\begin{continued}{tools.nex}
  -- What is wrong with the path the model gave, or "": it must stay in
  -- the workspace and, if `kind` is "file" or "folder", be one.
  path_problem(path, kind: String): String do
    result := ""
    if not workspace.contains(path) then
      result := "`" + path + "` is outside the workspace; use a path relative to it, such as `loans.nex`"
    elseif kind /= "any" and not workspace.exists(path) then
      result := "there is no " + kind + " `" + path + "`"
    elseif kind = "file" and not workspace.is_file(path) then
      result := "`" + path + "` is a folder, not a file"
    elseif kind = "folder" and not workspace.is_folder(path) then
      result := "`" + path + "` is a file, not a folder"
    end
  end
\end{continued}

Two private helpers finish the class. A parameter's specification is a
\code{Map[String, Any]} inside another one, so finding its type means
looking inside an \code{Any}. \code{convert} \ldots\ \code{to} does that
safely: \code{if convert x to m: Map[String, Any] then} runs its branch only
if \code{x} really is such a map, with \code{m} as that map.

\begin{continued}{tools.nex}
  private feature
    type_of_parameter(key: String): String do
      result := ""
      if convert parameters().get(key) to spec: Map[String, Any] then
        if convert spec.get("type") to json_type: String then
          result := json_type
        end
      end
    end

    is_whole_number(value: Json_Value): Boolean do
      result := false
      if convert value.raw() to n: Integer then
        result := true
      end
    end
end
\end{continued}

One more rule, for every tool: no result may be longer than about 5,000
tokens. A model can only read so much at once, and a whole log file sent
back by a tool would crowd out everything else.

\begin{append}{tools.nex}
-- The most text one tool result may hold (about 5,000 tokens). Longer
-- output is shortened, keeping its start and its end.
function output_limit(): Integer do
  result := 20000
end
\end{append}

\section{Listing files}

The first real tool lists what is in the workspace. It is the first thing
a programmer does in a project they don't know, and the description tells
the model to do the same. A heir's constructor sets up what it inherits by
calling its parent's, with \code{super.make}.

\begin{append}{tools.nex}
class List_Files inherit Tool
create
  make(ws: Workspace) do
    super.make(ws)
  end

feature
  name(): String do
    result := "list_files"
  end

  description(): String do
    result := "List the files and folders in the workspace, or in one folder of it, "
      + "including everything below it. Folders end in `/`. Use this first, to see "
      + "what the project contains."
  end

  parameters(): Map[String, Any] do
    result := {"path": parameter("string", "The folder to list, relative to the workspace. Default: `.`, the whole workspace.")}
  end

  required(): Array[String] do
    result := []
  end

  describe(input: Json_Value): String do
    result := "list_files " + string_argument(input, "path", ".")
  end

  run(input: Json_Value): Tool_Outcome do
    let path: String := string_argument(input, "path", ".")
    let problem: String := path_problem(path, "folder")
    if problem /= "" then
      result := create Tool_Outcome.error(problem)
    else
      let found: Array[String] := workspace.entries(path)
      if found.is_empty then
        result := create Tool_Outcome.ok("(the folder is empty)")
      elseif found.length > 300 then
        result := create Tool_Outcome.ok(joined_lines(found.take(300))
          + "... and " + (found.length - 300).to_string + " more")
      else
        result := create Tool_Outcome.ok(joined_lines(found))
      end
    end
  end
end
\end{append}

An empty folder gives \code{"(the folder is empty)"} rather than nothing at
all. A model that gets an empty result can't tell whether the tool found
nothing or failed.

\section{Reading files}

Reading a file numbers its lines. The numbers are for talking about the
file (``the bug is on line 32''); the description warns the model that they
are not part of the file, because it will later have to quote text from
it exactly. A long file is sent 500 lines at a time, and the result says
which lines it holds, so the model knows to ask for more.

\begin{append}{tools.nex}
class Read_File inherit Tool
create
  make(ws: Workspace) do
    super.make(ws)
  end

feature
  MOST_LINES = 500

  name(): String do
    result := "read_file"
  end

  description(): String do
    result := "Read a text file in the workspace. Lines come numbered, like `   12  code`; "
      + "the numbers are not part of the file. Read a file before you change it. "
      + "Long files come " + MOST_LINES.to_string + " lines at a time: use start_line to read on."
  end

  parameters(): Map[String, Any] do
    result := {"path": parameter("string", "The file, relative to the workspace."),
      "start_line": parameter("integer", "The first line to show, counting from 1. Default: 1."),
      "end_line": parameter("integer", "The last line to show. Default: the end of the file, or "
        + MOST_LINES.to_string + " lines on, whichever comes first.")}
  end

  required(): Array[String] do
    result := ["path"]
  end

  describe(input: Json_Value): String do
    result := "read_file " + string_argument(input, "path", "?")
    if integer_argument(input, "start_line", 1) > 1 then
      result := result + " from line " + integer_argument(input, "start_line", 1).to_string
    end
  end

  run(input: Json_Value): Tool_Outcome do
    let path: String := string_argument(input, "path", "")
    let problem: String := path_problem(path, "file")
    if problem /= "" then
      result := create Tool_Outcome.error(problem)
    else
      let lines: Array[String] := lines_of(workspace.read(path))
      let first: Integer := integer_argument(input, "start_line", 1).max(1)
      let last: Integer := integer_argument(input, "end_line", first + MOST_LINES - 1)
        .min(lines.length).min(first + MOST_LINES - 1)
      if lines.is_empty then
        result := create Tool_Outcome.ok("(the file is empty)")
      elseif first > lines.length then
        result := create Tool_Outcome.error("`" + path + "` has only "
          + count_text(lines.length, "line"))
      elseif last < first then
        result := create Tool_Outcome.error("end_line is before start_line")
      else
        let shown: String := numbered(lines.slice(first - 1, last), first)
        if first > 1 or last < lines.length then
          shown := shown + "(lines " + first.to_string + "-" + last.to_string
            + " of " + lines.length.to_string + ")\n"
        end
        result := create Tool_Outcome.ok(shortened(shown, output_limit()))
      end
    end
  end
end
\end{append}

\code{MOST_LINES = 500} is a \emph{constant}: a feature with a fixed value,
written in capitals by convention. \code{lines.slice(a, b)} is the lines
from position \code{a} up to, but not including, \code{b}, counting from 0;
line numbers count from 1, hence the \code{first - 1}.

Let's try both tools on a fresh copy of the sample project, with good input
and bad:

\begin{tryit}{try_tools.nex}
intern Tools

let ws: Workspace := fresh_workspace(user_path("sample_project"), user_path("scratch"))
let tools: Array[Tool] := [create List_Files.make(ws), create Read_File.make(ws)]

function try(tool: Tool, input: String) do
  let arguments: Json_Value := json_value(input)
  console.print_line("> " + tool.describe(arguments))
  let problem: String := tool.problem_with(arguments)
  if problem /= "" then
    console.print_line("bad input: " + problem)
  else
    let outcome: Tool_Outcome := tool.run(arguments)
    console.print_line((when outcome.is_error then "error: " else "" end) + outcome.text)
  end
end

try(tools.get(0), '{}')
try(tools.get(1), '{"path": "loans.nex", "start_line": 26, "end_line": 34}')
try(tools.get(1), '{"path": "../code.nex"}')
try(tools.get(1), '{"file": "loans.nex"}')
try(tools.get(1), '{"path": "loans.nex", "start_line": "26"}')
console.print_line(json_text(tools.get(0).definition()))
\end{tryit}
\tryoutput{try_tools}

The array of tools needs its type written out, \code{Array[Tool]}: its
elements are of two different classes, and Nex won't guess what they have
in common. In the next chapter, the agent will do what \code{try} does
here, for whatever tool the model asks for.
